\section{Розділ~\ref{sec:code}. Абетка Морзе}

Граничні оцінки розділу --- теорія інформації та рахунок; виміряно, де в каналі виникає шум, як швидко працює мозок і скільки коштує імпульс, а питання про те, яким кодом кора користується насправді, лишається відкритим.

{\footnotesize
\setlength{\tabcolsep}{3pt}\renewcommand{\arraystretch}{1.15}
\begin{longtable}{>{\raggedright}p{0.5cm}>{\raggedright}p{4.0cm}>{\raggedright}p{5.6cm}>{\raggedright}p{2.3cm}>{\raggedright\arraybackslash}p{1.7cm}}
\toprule
№ & Твердження & Дослід і що виміряно & Джерело & Статус \\
\midrule\endfirsthead
\toprule
№ & Твердження & Дослід і що виміряно & Джерело & Статус \\
\midrule\endhead
1 & Частоти \nt{GLUT}-нейронів кори --- від часток до десятків імпульсів за секунду, і більшу частину часу нейрони працюють біля нижньої межі & Запис через піпетку, прикладену до клітини (вибір нейрона не залежить від його активності), у слуховій корі щура, що не спить: у кожен момент частоту вище $20$ імпульсів за секунду дають менше $5\,\%$ нейронів; у частини нейронів фонова частота нижча за $0.01$ за секунду; розподіл частот логнормальний & \cite{Hromadka2008} & виміряно \\
2 & Пропускна здатність імпульсного каналу $R_{\max}\approx r\log_2\frac{e}{r\Delta t}$~\eqref{eq:spikeentropy}, майже вся вона --- у часі імпульсів (твердження~\ref{prop:capacity}) & Ентропія двійкового рядка з клітинок довжини $\Delta t$. Числа розділу: $8$ біт на імпульс при $r=10$ за секунду й $\Delta t=1\ms$, близько $5$ біт при $\Delta t=10\ms$, менше $2$ біт за секунду в лічильника імпульсів & \cite{MacKay1952} & теорія \\
3 & Чутливі нейрони передають від одного до кількох біт на імпульс, у найкращих випадках до половини межі & Відновлення подразника за імпульсами чутливих нейронів, для яких подразник відомий; зведення вимірювань у книзі & \cite{Rieke1997} & виміряно \\
4 & Генератор імпульсів точний; точний час задають швидкі перепади входу & Нейрони кори щура у зрізах: на повторюваний струм зі швидкими коливаннями імпульси повторюються з точністю менше $1\ms$, на сталий струм моменти розповзаються & \cite{Mainen1995} & виміряно \\
5 & Синапс ненадійний: понад половина імпульсів не доходить до отримувача через випадковість викиду & Мінімальна стимуляція входів пірамідних нейронів гіпокампа (зрізи): понад половина імпульсів не викликає відповіді. Виключено коливання порога стимуляції, збої проведення на розгалуженнях аксона й низьку температуру досліду; лишається ймовірнісний викид & \cite{AllenStevens1994} & виміряно \\
6 & Потік імпульсів у живій корі виглядає випадковим: інтервали розкидані як у пуассонівського потоку, дисперсія кількості імпульсів близька до середнього; частина «шуму» --- повідомлення про рухи тварини & Записи в зорових областях V1 і MT мавпи, що не спить: коефіцієнт варіації інтервалів близько $1$, відношення дисперсії кількості імпульсів до середнього --- як у випадкового процесу; модель, що підсумовує незалежні малі входи, дає розкид набагато менший. У зоровій корі миші значну частину розкиду передбачають за відеозаписом рухів самої тварини & \cite{Softky1993,Stringer2019} & виміряно \\
7 & Частотний код повільний: похибка $1/\sqrt{rT}$~\eqref{eq:rateerror}, а мозок упізнає картинку за $\sim150\ms$ & Формула --- пуассонівська статистика, висновок розділу. Дослід: люди вирішували, чи є тварина на незнайомій фотографії, показаній на $20\ms$; потенціали мозку на «так» і «ні» розходяться приблизно через $150\ms$. Розбиття цього часу на десяток ступенів по $10$--$20\ms$ --- оцінка розділу, не вимірювання & \cite{Thorpe1996} & виміряно \\
8 & Усереднення за групою обмежене кореляцією: похибка падає не більше ніж у $1/\sqrt c$ разів~\eqref{eq:correlated} (твердження~\ref{prop:pooling}) & Пари нейронів області MT мавпи, записані одночасно: коефіцієнт кореляції кількостей імпульсів у середньому $0.12$, і теоретичний аналіз показує, що навіть така кореляція обмежує виграш від групи. Теорія: межу ставить лише та частина кореляції, що схожа на сигнал. Модель протилежної позиції: частоту групи з $50$--$100$ нейронів читають за один міжімпульсний інтервал, $10$--$50\ms$, і час окремих імпульсів у корі не використовується & \cite{Zohary1994,MorenoBote2014,ShadlenNewsome1998} & виміряно \\
9 & Число можна писати в час першого імпульсу~\eqref{eq:latency}, у порядок імпульсів групи або у фазу місцевого ритму & Сітківка саламандри: просторову структуру коротко показаної картинки записано у відносних затримках перших імпульсів вихідних нейронів, майже незалежно від контрасту. Клітини місця гіпокампа щура: під час проходу через «своє» місце імпульси зсуваються дедалі раніше в циклі тета-ритму ($7$--$12$\,Гц), зсув від $100^\circ$ до $355^\circ$. Наскільки кора користується часом імпульсів --- відкрите питання (рядок~8) & \cite{Gollisch2008,OKeefe1993} & виміряно \\
10 & Який код прочитано, вирішує стала часу отримувача~\eqref{eq:reader} (твердження~\ref{prop:reader}); в активній корі нейрони ближчі до детекторів збігів & Формула~\eqref{eq:reader} --- висновок розділу. Огляд записів in vivo: в активній корі провідність мембрани в кілька разів вища, ніж у спокої, і стала часу відповідно коротша. Що кора перемикає код саме так, безпосередньо не показано & \cite{Destexhe2003} & гіпотеза \\
11 & Синапс, що виснажується, передає зміну частоти, по суті $\Delta r/r$~\eqref{eq:depression} (рис.~\ref{fig:depression}) & Парні записи пірамідних нейронів кори: швидкість виснаження задається ймовірністю викиду; за повільного виснаження відповідь відображає частоту, за швидкого --- збіги. Модель на основі виміряного виснаження в корі щура: однакові відносні зміни частоти на швидких і повільних входах дають однакову відповідь, тобто синапс передає $\Delta r/r$. Числа $1.7\to2.2$, $6.7$ і $90\ms$ --- розрахунок розділу & \cite{Tsodyks1997,Abbott1997depression} & виміряно \\
12 & Пачка --- «тире»: її ймовірність загубитися $(1-p)^k$~\eqref{eq:burst} мала, полегшення зменшує її ще & Огляд: на ненадійних синапсах одиночні імпульси часто губляться, а пачки передаються надійно завдяки полегшенню; пачка викликає довготривалі зміни синапсів. Що мозок читає пачку як окремий знак, --- гіпотеза & \cite{Lisman1997,AllenStevens1994} & гіпотеза \\
13 & Швидкі \nt{GABA}-нейрони задають ритм і «пунктуацію»; ще один клас гальмує гальмівних і відчиняє ворота (твердження~\ref{prop:punctuation}) & Огляд властивостей класів \nt{GABA}-нейронів кори. Оптогенетика: ритмічне збудження швидких \nt{GABA}-нейронів викликає гамма-ритм, і відповідь на подразник залежить від фази циклу. У зоровій корі миші PV-нейрони гальмують один одного, SST-нейрони --- усі інші класи, VIP-нейрони --- переважно SST. Збудження VIP-нейронів у миші, що не спить, знімає гальмування з \nt{GLUT}-нейронів; їх вмикає сигнал підкріплення. Поділ ролей як загальне правило --- гіпотеза & \cite{Tremblay2016,Cardin2009,Pfeffer2013,Pi2013} & гіпотеза \\
14 & Енергія обмежує середню частоту величиною порядку одного імпульсу за секунду~\eqref{eq:energy} & Енергетичний бюджет сірої речовини гризуна за анатомічними й фізіологічними даними: імпульси й синаптичні струми --- $47\,\%$ і $34\,\%$ енергії сигналізації; одночасно можуть бути активними не більше $\sim15\,\%$ нейронів. Для кори людини: не більше $\sim1\,\%$ нейронів можуть бути сильно активними одночасно. Число $\sim10^9$ АТФ на імпульс і потужність $\sim10$\,Вт на сигналізацію --- оцінка розділу на основі цих розрахунків & \cite{Attwell2001,Lennie2003} & теорія \\
15 & Код розріджений і підлаштований під найбільшу кількість біт на джоуль; кора міняє точність на енергію (твердження~\ref{prop:economy}) & Гіпотеза ощадливого коду. Опора: у мишей за обмеження їжі в зоровій корі падає провідність AMPA-рецепторів і синаптична витрата АТФ ($-29\,\%$), частота імпульсів зберігається, а налаштування на орієнтацію розширюється на $32\,\%$ & \cite{Barlow1961,Padamsey2022} & гіпотеза \\
\bottomrule
\end{longtable}}
